\chapter{Behind the Scenes of Diffusion Models:\\ It\^o’s Calculus and Girsanov's Theorem}\label{app:Ito}


(Score-based) Diffusion models are built on SDEs: a drift that pushes states and a Brownian term that jitters them. Unlike ODE paths, Brownian paths are nowhere differentiable, so the ordinary chain rule fails. In this section, we introduce two fundamental tools that make the math precise:
\begin{itemize}
\item \textbfs{It\^o’s Formula} is the correct chain rule for stochastic trajectories. It tells us how a function $\rvh(\rvx_t,t)$ evolves when $\rvx_t$ follows an SDE. It enables derivations of the Fokker–Planck equation, moment dynamics, the It\^o product rule, and the identities used in score-based training.
\item \textbfs{Girsanov’s Theorem} is a change-of-measure result on \emph{path probabilities}. It quantifies how likelihoods change when the noise is fixed but the drift is altered. This links score matching to path-space KL divergence and explains why learning the score in the reverse SDE corresponds to maximizing the data likelihood.
\end{itemize}

With these tools, the standard diffusion model derivations (Fokker–Planck, reverse time SDE, training objectives, and likelihood relations) follow cleanly and without hand waving.


\newpage


\section{It\^o's Formula: The Chain Rule for Random Processes}
Standard calculus does not directly apply to stochastic processes because Wiener processes are not differentiable in the classical sense. Instead, we use It\^o’s calculus, which provides rules for working with stochastic integrals. 

\subsection{Motivation: Why Do We Need a Special Chain Rule?}
Consider a deterministic time-varying function $\rvy_t$ that evolves smoothly with time $t$ (e.g., an ODE). If we have a function $\rvh(\rvy_t, t)$, the usual chain rule tells us:
\[
\frac{\diff \rvh}{\diff t} = \frac{\partial \rvh}{\partial t}  + \nabla_{\rvy}\rvh  \frac{\diff \rvy_t}{\diff t} .
\]
Here, $\nabla_{\rvy}\rvh$ is the Jacobian of $\rvh$. This works perfectly for deterministic paths $\rvy_t$.

\begin{question}
    But what happens if $\rvx_t$ is a stochastic process, say, driven by an SDE 
\[
\diff \rvx_t = \rvf(\rvx_t, t)\diff t + g(t) \diff \rvw_t
\]
as in \Cref{eq:sde}? What SDE does the process $\rvh(\rvx_t, t)$ satisfy?
\end{question}

\paragraph{Why the Ordinary Chain Rule Fails?}
Naïvely applying the classical chain rule yields
\[
\diff \rvh = \frac{\partial \rvh}{\partial t} \diff t + \nabla_{\rvy}\rvh  \cdot \diff \rvx_t.
\]
However, this neglects that Brownian increments satisfy $\diff \rvw_t = \mathcal{O}(\sqrt{\diff t})$ and
\[
(\diff \rvw_t)^2 = \diff t.
\]
Thus, second-order terms in $\diff \rvw_t$ \textit{do not} vanish in stochastic calculus, unlike classical calculus where $(\diff t)^2$ terms are negligible.

\exm{Simple Example--$h(x_t) = x_t^2$}{
To see the intuition, let us consider the simple real-valued function $h(x_t) = x_t^2 \in \mathbb{R}$ where the random variable $x_t \in \mathbb{R}$ satisfies
\[
\diff x_t = \sigma \diff w_t, 
\]
with a constant $\sigma>0$. 
If we try the classical chain rule,
\[
\diff h = 2 x_t \diff x_t = 2 x_t \sigma \diff w_t.
\]
If this were true, the expectation of $h(x_t)$ would be constant in time because 
\[
\mathbb{E}[\diff h ] =  2 \sigma \mathbb{E}[x_t  \diff w_t]= 2 \sigma \mathbb{E}[x_t]  \underbrace{\mathbb{E}[\diff w_t]}_{=0}=0.
\]
But we know from classical Brownian motion properties (see \Cref{eq:wiener-gaussian}) that
\[
\mathbb{E}[x_t^2] = \sigma^2 t,
\]
which grows linearly in time. So the ordinary chain rule misses an important term.
}


\subsection{Deriving 1D It\^o's Formula from Taylor Expansion}\label{subsec:1D-Ito}

\paragraph{Deterministic Chain Rule via Taylor Expansion.}
To understand why the classical chain rule fails for stochastic processes defined by SDEs, we first revisit it in the deterministic setting using Taylor expansion. We consider the scalar case: $y_t \in \mathbb{R}$ and $h(\cdot,\cdot)\in \mathbb{R}$. Formally treating $\diff y_t = y_{t+\diff t} - y_t$, with $\diff t \approx 0$, we expand:
\begin{align*}
&h(y_{t+\diff t}, t+\diff t) 
- h(y_t, t) \\
= ~&\frac{\partial h}{\partial t} \diff t
+ \frac{\partial h}{\partial y} \diff y_t  + \frac{1}{2} \left(\textcolor{gray}{\frac{\partial^2 h}{\partial y^2}  (\diff y_t)^2}
+ 2 \textcolor{gray}{\frac{\partial^2 h}{\partial t \partial y} \diff t \diff y_t}
+  \textcolor{gray}{\frac{\partial^2 h}{\partial t^2} (\diff t)^2}
\right)
+ \mathcal{O}(\diff t^3),
\end{align*}
Here, $\diff t \diff y_t = \left(\frac{\diff y_t}{\diff t}\right)(\diff t)^2 = \mathcal{O}(\diff t^2)$, and similarly $(\diff t)^2 = \mathcal{O}(\diff t^2)$. Therefore, all the gray parts are ignorable, and the full differential is:
\[
\diff h
= \frac{\partial h}{\partial t} \diff t
+ \frac{\partial h}{\partial y} \diff y_t
+ \mathcal{O}(\diff t^2).
\]
% In the limit $\diff t \to 0$, we recover the classical chain rule.

\paragraph{It\^o's Formula via Stochastic Taylor Expansion.}
Now consider a stochastic process $x_t \in \mathbb{R}$ governed by the SDE:
\[
\diff x_t = f(x_t, t) \diff t + g(t) \diff w_t,
\]
where $w_t$ is standard Brownian motion. We aim to compute the differential of a scalar-valued function $h(x_t, t)$.

Using the stochastic Taylor expansion~\citep{kloeden1992stochastic}, which retains second-order terms in $\diff x_t$, we have:
\begin{align*}
&h(x_{t+\diff t}, t+\diff t) 
- h(x_t, t) \\
=~&\frac{\partial h}{\partial t} \diff t 
+ \frac{\partial h}{\partial x} \diff x_t 
+ \frac{1}{2} \left( \frac{\partial^2 h}{\partial x^2} (\diff x_t)^2 
+ 2\textcolor{gray}{{\frac{\partial^2 h}{\partial t \partial x} \diff t \diff x_t}} 
+ \textcolor{gray}{{ \frac{\partial^2 h}{\partial t^2} (\diff t)^2}} \right)
+ \cdots
\end{align*}
\subparagraph{Negligible Cross Terms.}
By the scaling property of Brownian motion (\Cref{eq:wiener-gaussian}), 
\[
\diff w_t = \mathcal{O}(\sqrt{\diff t}) \quad \Rightarrow \quad \diff t \cdot \diff w_t = \mathcal{O}((\diff t)^{3/2}).
\]
Therefore,
\begin{align*}
\diff t \cdot \diff x_t 
&= \diff t \left(f \diff t + g \diff w_t\right) 
= f (\diff t)^2 + g \cdot \diff t \cdot \diff w_t 
= \mathcal{O}((\diff t)^{3/2}).
\end{align*}
So, the gray terms are negligible: $ \mathcal{O}((\diff t)^{3/2}) $ or smaller.
\subparagraph{Second-Order Term $(\diff x_t)^2$.}
Expanding using the SDE:
\begin{align*}
(\diff x_t)^2 
&= \left( f \diff t + g \diff w_t \right)^2 \\
&= f^2 (\diff t)^2 + 2fg \diff t \diff w_t + g^2 (\diff w_t)^2 \\
&= \mathcal{O}((\diff t)^2) + \mathcal{O}((\diff t)^{3/2}) + g^2 \mathcal{O}(\diff t) \\
&= g^2(t) \diff t + \mathcal{O}((\diff t)^{3/2}).
\end{align*}

Combining terms, we obtain the differential:
\[
\diff h(x_t, t)
= \frac{\partial h}{\partial t} \diff t
+ \frac{\partial h}{\partial x} \diff x_t
+ \frac{1}{2} \frac{\partial^2 h}{\partial x^2} g^2(t) \diff t.
\]
Substituting $ \diff x_t = f(x_t, t) \diff t + g(t) \diff w_t $ yields:
\[
\diff h(x_t, t)
= \left( \frac{\partial h}{\partial t} + f \frac{\partial h}{\partial x}
+ \frac{1}{2} g^2 \frac{\partial^2 h}{\partial x^2} \right) \diff t
+ g \frac{\partial h}{\partial x} \diff w_t.
\]
This is the 1D version of \emph{It\^o's formula}.


\exm{Simple Example--$h(x_t) = x_t^2$}{
We revisit the simple example: $ h(x_t) = x_t^2 $, where the stochastic process $ x_t \in \mathbb{R} $ satisfies
\[
\diff x_t = \sigma \diff w_t,
\]
with a constant $ \sigma > 0 $.  
 Applying It\^o’s formula correctly to $ h(x_t) = x_t^2 $, we obtain:
\[
\diff h(x_t) = \diff (x_t^2) = 2 x_t \diff x_t + \sigma^2 \diff t.
\]
Substituting $ \diff x_t = \sigma \diff w_t $, this becomes:
\[
\diff(x_t^2) = 2 x_t \sigma \diff w_t + \sigma^2 \diff t.
\]
}

\subsection{It\^o's Formula: The Chain Rule for SDEs}
\label{subsec:ito-summary}
We summarize the one-dimensional It\^o's formula derived above. Using similar arguments, the result extends naturally to the multi-dimensional setting. While we omit the detailed derivation, we state the general formula for completeness.

Finally, we illustrate an application of It\^o's formula by deriving the \emph{It\^o product rule}, which enables computation of $\diff(\rvx_t^\top \rvy_t)$ for stochastic processes $\rvx_t$ and $\rvy_t$.

\paragraph{$1$D It\^o's Formula.}
Let $x_t \in \mathbb{R}$ be a stochastic process satisfying the SDE:
\[
\diff x_t = f(x_t, t) \diff t + g(t) \diff w_t.
\]
For a scalar function $h \colon \mathbb{R} \times [0, T] \to \mathbb{R}$, the process $h(x_t, t)$ satisfies:
\begin{mdframed}
\[
\diff h(x_t, t)
= \left( \frac{\partial h}{\partial t} + f \frac{\partial h}{\partial x}
+ \frac{1}{2} g^2 \frac{\partial^2 h}{\partial x^2} \right) \diff t
+ g \frac{\partial h}{\partial x} \diff w_t.
\]
\end{mdframed}


\paragraph{Multidimensional It\^o's Formula with Scalar Output.}
Let $\rvx_t \in \mathbb{R}^D$ satisfy the SDE:
\[
\diff \rvx_t = \rvf(\rvx_t, t) \diff t + g(t) \diff \rvw_t,
\]
where $\rvf \colon \mathbb{R}^D \times [0, T] \to \mathbb{R}^D$, $g \colon [0, T] \to \mathbb{R}$, and $\rvw_t \in \mathbb{R}^D$ is a $D$-dimensional Brownian motion. Let $h \colon \mathbb{R}^D \times [0, T] \to \mathbb{R}$ be a scalar-valued function. Then $h(\rvx_t, t)$ satisfies:
\begin{mdframed}
\begin{align}\label{eq:Ito-scalar}
    \diff h(\rvx_t, t)
= \left( \frac{\partial h}{\partial t}
+ \nabla_{\rvx} h^\top \rvf
+ \frac{1}{2} g^2(t) \operatorname{Tr} \left( \nabla^2_{\rvx} h \right) \right) \diff t
+ g(t) \nabla_{\rvx} h^\top \diff \rvw_t,
\end{align}
\end{mdframed}
where $\nabla_{\rvx} h \in \mathbb{R}^D$ is the gradient and $\nabla^2_{\rvx} h \in \mathbb{R}^{D \times D}$ is the Hessian matrix of $h$ with respect to $\rvx$.


\exm{It\^o's Product Rule}{
Let $\rvx_t, \rvy_t \in \mathbb{R}^D$ be vector-valued stochastic processes governed by the SDEs:
\[
\begin{aligned}
\diff \rvx_t &= \rva(\rvx_t, t) \diff t + b(t) \diff \rvw_t, \\
\diff \rvy_t &= \rvc(\rvy_t, t) \diff t + d(t) \diff \rvw_t,
\end{aligned}
\]
where $\rva, \rvc: \mathbb{R}^D \times [0, T] \to \mathbb{R}^D$ are vector fields, and $b(t), d(t) \in \mathbb{R}$ are scalar-valued functions. Here, $\rvw_t \in \mathbb{R}^D$ denotes a standard $D$-dimensional Brownian motion.

We aim to derive the SDE for the scalar-valued process
\[
z(t) := \rvx_t^\top \rvy_t.
\]

Applying the multivariate It\^o formula to the bilinear function $h(\rvx, \rvy) := \rvx^\top \rvy$, we obtain:
\[
\diff(\rvx^\top \rvy) = (\diff \rvx)^\top \rvy + \rvx^\top \diff \rvy + \Tr\left[ \diff \rvx \cdot (\diff \rvy)^\top \right].
\]

The It\^o correction term is computed as:
\begin{align*}
    \diff \rvx \cdot (\diff \rvy)^\top &= b(t) \diff \rvw_t \cdot \left[ d(t) \diff \rvw_t \right]^\top \\&= b(t) d(t) \diff \rvw_t \cdot \diff \rvw_t^\top \\&= b(t)d(t) \diff t \cdot \rmI_D.
\end{align*}
Thus,
\[
\Tr\left[ \diff \rvx \cdot (\diff \rvy)^\top \right] = b(t)d(t) \Tr(\rmI_D) \diff t = D  b(t)d(t) \diff t.
\]

Putting everything together, the resulting SDE is:
\begin{align}\label{eq:ito-product}
\begin{aligned}
        \diff(\rvx^\top \rvy) &= (\diff \rvx)^\top \rvy + \rvx^\top \diff \rvy + D  b(t)d(t) \diff t
    % \\ &=\rva^\top \rvy \diff t + b(t) \rvy^\top \diff \rvw_t + \rvx^\top \rvc \diff t + d(t) \rvx^\top \diff \rvw_t + D  b(t)d(t) \diff t.
\end{aligned}
\end{align}
}



% \paragraph{Multidimensional It\^o's Formula with Vector Output.} Now, if  $\rvh \colon \mathbb{R}^D \times [0, T] \to \mathbb{R}^D$ is vector-valued. Then $\rvh(\rvx_t, t)$ satisfies:
% \begin{mdframed}
% \begin{align}\label{eq:Ito-vector}
%     \diff \rvh(\rvx_t, t)
% = \left(
% \frac{\partial \rvh}{\partial t}
% + \nabla_{\rvx}\rvh \cdot \rvf
% + \frac{1}{2} g^2(t) \sum_{i=1}^D \frac{\partial^2 \rvh}{\partial x_i^2}
% \right) \diff t
% + g(t) \nabla_{\rvx}\rvh \cdot\diff \rvw_t,
% \end{align}
% \end{mdframed}
% where $\nabla_{\rvx}\rvh \in \mathbb{R}^{D \times D}$ is the Jacobian matrix of $\rvh$ with respect to $\rvx$, and $\frac{\partial^2 \rvh}{\partial x_i^2} \in \mathbb{R}^D$ is the vector of second derivatives of each component of $\rvh$ with respect to $x_i$.

% Equivalently, in component-wise form: if $\rvh = (h^{(1)}, \dots, h^{(D)})^\top$, then for each $k = 1, \dots, D$,
% \begin{align*}
%     &\diff h^{(k)}(\rvx_t, t)
% \\= &\left(
% \frac{\partial h^{(k)}}{\partial t}
% + \nabla_{\rvx} h^{(k)} \cdot \rvf
% + \frac{1}{2} g^2(t) \operatorname{Tr} \left( \nabla_{\rvx}^2 h^{(k)} \right)
% \right) \diff t
% + g(t) \nabla_{\rvx} h^{(k)}  \diff \rvw_t.
% \end{align*}




\subsection{It\^o's Formula's Application: Derivation of Fokker-Planck Equation}\label{subsec:rigorous-proof-fpe}
In this section, we apply It\^o's formula from \Cref{eq:Ito-scalar} to derive the Fokker–Planck equation, a PDE that characterizes the time evolution of the probability density $ p_t(\rvx) $ associated with the $ D $-dimensional diffusion process defined by the SDE in \Cref{eq:sde}.

\paragraph{Step 1: Apply It\^o's Formula.}
Let $\phi(\rvx, t)$ be a smooth test function $\phi \colon \mathbb{R}^D \times [0, T] \to \mathbb{R}$. By It\^o’s Formula:
\[
\diff \phi(\rvx_t, t) = \left( \frac{\partial \phi}{\partial t} + \nabla_{\rvx} \phi^\top \rvf(\rvx_t, t) + \frac{1}{2} g^2(t) \Tr[\nabla_{\rvx}^2 \phi] \right) \diff t + g(t) \nabla_{\rvx} \phi^\top \diff \rvw_t.
\]
\paragraph{Step 2: Take Expectation.}
Taking expectation over $p_t(\rvx)$ and noting $\mathbb{E}[\diff \rvw_t] = 0$:
\[
\mathbb{E}[\diff \phi(\rvx_t, t)] = \mathbb{E} \left[ \left( \frac{\partial \phi}{\partial t} + \nabla_{\rvx} \phi^\top \rvf(\rvx_t, t) + \frac{1}{2} g^2(t) \Tr[\nabla_{\rvx}^2 \phi] \right) \diff t \right].
\]
\paragraph{Step 3: Express Expectation via Density.}
This expectation can be written as:
\[
\mathbb{E}[\diff \phi(\rvx_t, t)] = \int \left( \frac{\partial \phi}{\partial t} + \nabla_{\rvx} \phi^\top \rvf(\rvx, t) + \frac{1}{2} g^2(t) \Tr[\nabla_{\rvx}^2 \phi] \right) p_t(\rvx) \diff \rvx \diff t.
\]
\paragraph{Step 4: Integrate by Parts.}
Use integration by parts (divergence theorem) in $\mathbb{R}^D$:
\begin{align*}
    \int \nabla_{\rvx} \phi^\top \rvf p_t \diff \rvx &= -\int \phi \nabla_{\rvx} \cdot (\rvf p_t) \diff \rvx, \\
    \int \Tr[\nabla_{\rvx}^2 \phi] p_t  \diff \rvx &= \int \phi \Delta p_t \diff \rvx.
\end{align*}
\paragraph{Step 5: Substitute and Rearrange.}
Note that
\[
\mathbb{E}[\phi(\rvx_t,t)] = \int \phi(\rvx,t) p_t(\rvx) \diff\rvx.
\]
Differentiating in time, we have
\[
\frac{\diff}{\diff t}\mathbb{E}[\phi(\rvx_t,t)]
= \frac{\diff}{\diff t}\int \phi(\rvx,t) p_t(\rvx) \diff\rvx
= \int\Bigl(\partial_t\phi(\rvx,t) p_t(\rvx) + \phi(\rvx,t) \partial_t p_t(\rvx)\Bigr)\diff\rvx.
\]
Therefore,
\[
\mathbb{E}[\diff \phi(\rvx_t,t)]
= \frac{\diff}{\diff t}\mathbb{E}[\phi(\rvx_t,t)] \diff t
= \int\Bigl(\partial_t\phi(\rvx,t) p_t(\rvx) + \phi(\rvx,t) \partial_t p_t(\rvx)\Bigr)\diff\rvx \diff t.
\]
Equating this with Step~3, the terms $\int \partial_t\phi\,p_t\,\diff \rvx\,\diff t$ cancel. Applying Step~4 to the remaining spatial terms yields
\[
\mathbb{E}[\diff \phi(\rvx_t, t)]
= \int \phi(\rvx, t)\left[ \frac{\partial p_t}{\partial t}
+ \nabla_{\rvx}\cdot(\rvf(\rvx,t)\,p_t(\rvx))
- \frac{1}{2} g^2(t)\,\Delta p_t(\rvx) \right]\diff \rvx\,\diff t.
\]

\subsection{It\^o's Formula Application: Closed-Form Solution of a Linear SDE}\label{subsec:rigorous-proof-linear-sde}
This subsection demonstrates how to obtain a closed-form solution for a linear SDE by using an integration factor (similar to the ODE case) and It\^o's formula. The approach mirrors classical techniques for solving linear ODEs, but adapted to the stochastic setting.

We consider a linear SDE of the form
\begin{align}\label{eq:app-linear-sde}
    \mathrm{d}\rvx_t = f(t)\rvx_t \mathrm{d}t + g(t) \mathrm{d}\rvw_t,
\end{align}
where $f(t)$ and $g(t)$ are deterministic functions and $\rvw_t$ is a standard Wiener process. 
\paragraph{Closed-Form Solution of a Linear SDE.} We derive the explicit solution to the linear SDE in \Cref{eq:app-linear-sde} using the method of integrating factors. This type of forward SDE commonly arises in diffusion models (see \Cref{sec:dm-today}).


\subparagraph{Step 1: Define an Integration Factor.} Let
\[
\Psi(t) := \exp\left( -\int_0^t f(s) \mathrm{d}s \right),
\quad \text{and define} \quad
\rvy_t := \Psi(t)\rvx_t.
\]
\subparagraph{Step 2: Apply It\^o's Formula.}  
We apply It\^o's formula to the function $\rvh(\rvx, t) := \Psi(t)\rvx$. This is actually a special case of the It\^o product rule in \Cref{eq:ito-product}. Since $\Psi(t)$ is deterministic, there is no cross-variation term, and the formula simplifies to:
\begin{align*}
    \mathrm{d}\rvy_t 
    &= \mathrm{d}[\Psi(t)\rvx_t] \\
    &= \Psi'(t)\rvx_t \mathrm{d}t + \Psi(t) \mathrm{d}\rvx_t \\
    &= -f(t)\Psi(t)\rvx_t \mathrm{d}t + \Psi(t)\left[f(t)\rvx_t \mathrm{d}t + g(t) \mathrm{d}\rvw_t\right] \\
    &= \Psi(t)g(t) \mathrm{d}\rvw_t.
\end{align*}
Hence,
\[
\rvy_t = \rvy_0 + \int_0^t \Psi(s)g(s) \mathrm{d}\rvw(s) = \rvx_0 + \int_0^t \Psi(s)g(s) \mathrm{d}\rvw(s),
\]
since $\Psi(0) = 1$.
\subparagraph{Step 3: Solve for $\rvx_t$.}  
Using $\rvx_t = \Psi(t)^{-1}\rvy_t$, we obtain
\begin{align}\label{eq:lienar-sde-solution}
    \rvx_t &=  e^{ \int_0^t f(s)\mathrm{d}s}
    \left[ \rvx_0 + \int_0^t e^{ -\int_0^s f(r) \mathrm{d}r } g(s) \mathrm{d}\rvw(s) \right].
\end{align}
This provides an explicit solution to the vector-valued SDE.

Below, we demonstrate two alternative approaches to compute the analytical form of $p_t(\rvx_t|\rvx_0)$.

\paragraph{Analysis of the Closed-Form Solution.} \Cref{eq:lienar-sde-solution} reconfirms that $p_t(\rvx_t|\rvx_0)$ is Gaussian. To see this, define
\[
\phi(s) := e^{-\int_0^s f(u) \diff u} g(s),
\]
which is a deterministic matrix-valued function of time (assuming $f(u)$ and $g(s)$ are deterministic). The It\^o integral $\int_0^t \phi(s) \diff \rvw_s$ is then a zero-mean Gaussian random variable, as it is the stochastic integral of a deterministic function with respect to Brownian motion. Therefore, $\rvx_t|\rvx_0$ is an affine transformation of a Gaussian random variable and hence itself Gaussian. Its distribution is fully characterized by its conditional mean and covariance.

We define the conditional mean and covariance (given initial condition $\rvx_0$) as
\[
\rvm(t) := \mathbb{E}[\rvx_t | \rvx_0], \quad \rmP(t) := \mathbb{E}[(\rvx_t - \rvm(t))(\rvx_t - \rvm(t))^\top|\rvx_0].
\]

\subparagraph{Mean.}
Using linearity of expectation and the fact that the It\^o integral of a deterministic function has zero mean:
\begin{align*}
    \rvm(t) = \mathbb{E}\left[ e^{\int_0^t f(s) \diff s} \left( \rvx_0 + \int_0^t \phi(s) \diff \rvw_s \right) \bigg| \rvx_0 \right] = e^{\int_0^t f(s) \diff s} \rvx_0.
\end{align*}

\subparagraph{Covariance.}
Let  $\rvz_t := \int_0^t \phi(s) \diff \rvw_s$. Then $\rvx_t - \rvm(t) = A(t) \rvz_t$, so
\[
\rmP(t) = e^{2\int_0^t f(s) \diff s} \mathbb{E}[\rvz_t \rvz_t^\top].
\]
By It\^o isometry\footnote{It\^o's isometry links stochastic integrals to standard integrals in expectation; we omit the proof as it requires the full machinery of It\^o calculus. For a process $\bm{\psi} : [0,T] \to \mathbb{R}^{D \times D}$, It\^o isometry states
\[
\mathbb{E}\left[\left\| \int_0^T \bm{\psi}(t)   \mathrm{d}\rvw_t \right\|^2 \right]
= \mathbb{E}\left[\int_0^T \|\bm{\psi}(t)\|_F^2   \mathrm{d}t \right],
\]
where $\|\bm{\psi}(t)\|_F^2 = \sum_{i,j=1}^D |\psi_{ij}(t)|^2$ is the Frobenius norm. For $\bm{\psi}(t) \in \mathbb{R}^D$ (a vector), the integral is scalar and the isometry simplifies to
\[
\mathbb{E}\left[\left(\int_0^T \bm{\psi}(t)   \mathrm{d}\rvw_t\right)^2\right] = \mathbb{E}\left[\int_0^T \|\bm{\psi}(t)\|^2   \mathrm{d}t \right].
\]
},
\[
\mathbb{E}[\rvz_t \rvz_t^\top] = \left( \int_0^t \phi^2(s) \diff s \right) \rmI_D,
\]
hence,
\[
\rmP(t) = e^{2\int_0^t f(s) \diff s} \left( \int_0^t \left( e^{-\int_0^s f(u) \diff u} g(s) \right)^2 \diff s \right) \rmI_D.
\]
This shows the conditional covariance is isotropic.

\paragraph{Derivation of Mean and Variance ODEs in \Cref{eq:mean-var-evolution}.}
Alternatively, we can derive the moment evolution equations directly from the linear SDE \Cref{eq:app-linear-sde}. 


\subparagraph{Mean Evolution.}
Taking the conditional expectation of both sides of the SDE and using linearity:
\[
\frac{\diff \rvm(t)}{\diff t} = \mathbb{E}[f(t) \rvx_t|\rvx_0] = f(t) \mathbb{E}[\rvx_t|\rvx_0] = f(t) \rvm(t).
\]

\vspace{0.3cm}

\subparagraph{Covariance Evolution.}
Define the centered process $ \tilde{\rvx}_t := \rvx_t - \rvm(t) $. Applying It\^o's product rule (see \Cref{eq:ito-product}):
\[
\diff \left( \tilde{\rvx}_t \tilde{\rvx}_t^\top \right)
= \diff \tilde{\rvx}_t \cdot \tilde{\rvx}_t^\top 
+ \tilde{\rvx}_t \cdot \diff \tilde{\rvx}_t^\top 
+ \diff \tilde{\rvx}_t \cdot \diff \tilde{\rvx}_t^\top.
\]

From the SDE, we compute:
\[
\diff \tilde{\rvx}_t = \diff \rvx_t - \diff \rvm(t) = f(t) \tilde{\rvx}_t  \diff t + g(t)  \diff \rvw_t.
\]

Substituting into the product rule and taking expectation:
\begin{align*}
    \frac{\diff \rmP(t)}{\diff t}
    &= \mathbb{E}[f(t) \tilde{\rvx}_t \tilde{\rvx}_t^\top + \tilde{\rvx}_t \tilde{\rvx}_t^\top f(t) + g^2(t) \rmI_D] \\
    &= 2 f(t) \rmP(t) + g^2(t) \rmI_D.
\end{align*}
Thus, we recover the moment evolution equations in \Cref{eq:mean-var-evolution}.




% \subsection{Reverse-Time SDE Obeys the Same Fokker–Planck Equation as the Forward SDE}\label{subsec:reverse-sde-fokker-planck} 

% As discussed in \Cref{subsec:fpe-sde-ode}, the reverse-time SDE is constructed to match the marginal distributions of the forward SDE, thereby enabling the transformation of the noise distribution back into the data distribution. In this section, we formally derive the reverse-time SDE (see \Cref{subsec:reverse-sde-reason}) and verify that it satisfies the same Fokker-Planck equation as the forward process, ensuring consistency in marginal dynamics.

% Let $p_t(\rvx)$ denote the marginal distribution of the stochastic process $\rvx(t)$ defined by the forward SDE in \Cref{eq:sde}:
% \[
%     \diff \rvx(t) = \rvf(\rvx(t), t)\diff t + g(t) \diff \rvw(t),
% \]
% from $t=0$ to $t=T$. The corresponding Fokker–Planck equation that governs its marginal density $p_t(\rvx)$  is
% \begin{equation} \label{eq:forward_fp}
% \frac{\partial p_t}{\partial t} = -\nabla_{\rvx} \cdot \left( \rvf(\rvx, t) p_t(\rvx) \right) + \frac{1}{2} g^2(t) \Delta p_t(\rvx).
% \end{equation}

 




% \paragraph{Goal.} We aim to derive the reverse-time SDE that describes the dynamics of the time-reversed process. Specifically, we seek an equation of the form
% \begin{equation} \label{eq:reverse_sde_form}
% \diff \bar{\rvx}(t) = \bar{\rvf}(\bar{\rvx}(t), t) \diff t + \bar{g}(t) \diff \bar{\rvw}(t),
% \end{equation}
% where $\bar{\rvw}(t)$ denotes a Brownian motion adapted to reverse time, and $\bar{g}(t)$ is the diffusion coefficient. We require that the marginals of the reverse process $\bar{\rvx}(t)$ match those of the original forward process at reversed time, that is,
% \[
% \bar{p}_t(\rvx) := p_{T - t}(\rvx).
% \]
% In other words, the time-reversed process should reproduce the same family of distributions as the forward process, just traced backward in time. This guarantees that both forward and reverse processes follow consistent Fokker–Planck dynamics for their respective directions of evolution.



% \paragraph{Derivation of Reverse-Time SDE That Shares the Same Fokker–Planck Equation.}
%  Since $\bar p_t(\mathbf{x})=p_{T-t}(\mathbf{x})$, differentiating with respect to $t$ (and letting $s = T-t$, so $\diff s = -\diff  t$) gives
% \[
% \frac{\partial\bar p_t}{\partial t} = \frac{\partial p_{T-t}}{\partial (T-t)} \frac{\partial (T-t)}{\partial t} = - \frac{\partial p_s}{\partial s}\Big|_{s=T-t}.
% \]
% Substituting the Fokker–Planck \Cref{eq:forward_fp} for $p_s$ (with $s=T-t$), we get
% \begin{align*}
%     \frac{\partial\bar p_t}{\partial t} &= -\left[ -\nabla_{\mathbf{x}}\!\cdot\bigl(\rvf(\mathbf{x}, T-t) \bar p_t(\mathbf{x})\bigr) + \frac{1}{2} g^2(T-t) \Delta \bar p_t(\mathbf{x}) \right]
%  \\ &= \nabla_{\mathbf{x}}\!\cdot\bigl(\rvf(\mathbf{x}, T-t) \bar p_t(\mathbf{x})\bigr) - \frac{1}{2} g^2(T-t) \Delta \bar p_t(\mathbf{x})
%  \\ &= \left[\nabla_{\mathbf{x}}\!\cdot\bigl(\rvf(\mathbf{x}, T-t) \bar p_t(\mathbf{x}) \bigr) - g^2(T-t) \Delta \bar p_t(\mathbf{x}) \right] + \frac{1}{2} g^2(T-t) \Delta \bar p_t(\mathbf{x}).
% \end{align*}
% On the other hand, plugging \Cref{eq:reverse_sde_form} into its forward Fokker–Planck yields
% \[
% \frac{\partial\bar p_t}{\partial t} = - \nabla_{\mathbf{x}}\!\cdot\bigl(\bar \rvf(\mathbf{x}, t) \bar p_t(\mathbf{x})\bigr) + \tfrac12 \bar g^2(t) \Delta \bar p_t(\mathbf{x}).
% \]
% Matching diffusion terms gives $\bar g(t)=g(T-t)$. By matching drift terms (which is rigorously justified by Girsanov’s theorem but here derived through a formal calculation), we obtain
% \[
% -\nabla_{\mathbf{x}}\!\cdot(\bar \rvf  \bar p) = \nabla_{\mathbf{x}}\!\cdot\bigl(\rvf  \bar p\bigr) - g^2 \Delta \bar p = \nabla_{\mathbf{x}}\!\cdot\bigl(\rvf  \bar p - g^2 \nabla_{\mathbf{x}}\bar p\bigr).
% \]
% Hence,
% \[
% -\bar \rvf(\mathbf{x}, t)  \bar p_t = \rvf(\mathbf{x}, T-t)  \bar p_t - g^2(T-t) \nabla_{\mathbf{x}}\bar p_t,
% \]
% so
% \begin{align*}
%     \bar \rvf(\mathbf{x}, t) &= -\rvf(\mathbf{x}, T-t) + g^2(T-t) \frac{\nabla_{\mathbf{x}}\bar p_t(\mathbf{x})}{\bar p_t(\mathbf{x})} \\&= -\rvf(\mathbf{x}, T-t) + g^2(T-t) \nabla_{\mathbf{x}}\log\bar p_t(\mathbf{x}).
% \end{align*}
% Since $\bar p_t(\mathbf{x})=p_{T-t}(\mathbf{x})$, we conclude
% \[
% \mathrm{d} \bar{\mathbf{x}}(t) = \Bigl[-\rvf(\bar{\mathbf{x}}(t), T-t) + g^2(T-t) \nabla_{\mathbf{x}}\log p_{T-t}(\bar{\mathbf{x}}(t))\Bigr] \mathrm{d}t + g(T-t) \mathrm{d} \overline \rvw(t).
% \]

% \paragraph{Equivalent Reverse‐Flow Form.} Setting $t' = T-t$ (so $\diff t = -\diff t'$), and noting $\bar{\mathbf{x}}(t) = \mathbf{x}(T-t) = \mathbf{x}(t')$, we can rewrite the SDE in terms of the original time‐axis $t'$ running from $T$ down to $0$:
% \[
% \mathrm{d} \mathbf{x}(t') = \Bigl[-\rvf(\mathbf{x}(t'), t') + g^2(t') \nabla_{\mathbf{x}}\log p_{t'}(\mathbf{x}(t'))\Bigr] (-\mathrm{d}t') + g(t') \mathrm{d} \overline \rvw(T-t').
% \]
% Renaming $t'$ to $t$ for notational convenience (so that $t$ now runs backward from $T$ to $0$), and viewing $\mathrm{d}\overline{\rvw}(T-t)$ as a new reverse-time Brownian motion (still denoted by $\mathrm{d}\overline{\rvw}(t)$), we obtain:
% \[
% \mathrm{d} \mathbf{x}(t) = \bigl[\rvf(\mathbf{x}(t), t) - g^2(t) \nabla_{\mathbf{x}} \log p_t(\mathbf{x}(t))\bigr] \mathrm{d}t + g(t) \mathrm{d} \overline{\rvw}(t),
% \]
% with $t$ decreasing from $T$ to $0$.



\newpage

\section{Change-of-Variable For Measures: Girsanov's Theorem in Diffusion Models}
Diffusion models harness SDEs to transform simple noise into rich data distributions. At the heart of this transformation lies a profound idea: we can reinterpret randomness by modifying only the deterministic part of an SDE (the drift) while preserving its underlying stochasticity. This is precisely where Girsanov’s theorem enters the picture.


\paragraph{The Core Idea.}
Consider an observed continuous trajectory that describes the data's evolution from time $t=0$ to $t=T$, denoted as $\rvx_{0:T}:=\{ \rvx_t|t \in [0,T] \}$. Girsanov's theorem addresses a fundamental question:
\begin{question}
    Given this single observed path, what is its likelihood if we assume it was generated by one SDE, versus if we assume it was generated by a different SDE?
\end{question}
We compare two hypothetical models for generating the same trajectory. Both of these assumed SDEs share the same underlying pure randomness, represented by a standard Wiener process (Brownian motion) $\rvw_t$, but differ only in their deterministic ``push'' or ``drift'' function. We assume $\rvx_0$ has the same initial distribution for both assumed generating processes.

To build intuition, imagine $\rvx_{0:T}$ as a wiggly line drawn on paper. One hypothesis is that it was produced by a ``robot painter'' guided by a drift $\rvf$ and perturbed by random noise scaled by $g(t)$, yielding likelihood $p_{\rvf}(\rvx_{0:T})$. Alternatively, we imagine a second robot, with a different drift $\tilde{\rvf}$ but using the same noise process, generating the same line with likelihood $p_{\tilde{\rvf}}(\rvx_{0:T})$. Girsanov’s theorem gives us a precise way to compare these two likelihoods for the exact same observed path. It quantifies how a change in drift affects the probability of generating a particular trajectory, while holding the randomness fixed.

\paragraph{The Setup.}
Let $\rvx_t \in \mathbb{R}^D$ be our single, fixed, continuous path. We consider its likelihood under two SDE models, which differ only in their drift functions $\rvf$ and $\tilde{\rvf}$. They share the same diffusion coefficient $g(t) \in \mathbb{R}$ and the same underlying Wiener process $\rvw_t$:
\begin{align*}
    \diff\rvx_t &= \rvf(\rvx_t, t)\diff t + g(t)\diff\rvw_t \quad (\text{Model with drift } \rvf) \\
    \diff\rvx_t &= \tilde{\rvf}(\rvx_t, t)\diff t + g(t)\diff\rvw_t \quad (\text{Model with drift } \tilde{\rvf})
\end{align*}
Let $\bm{\delta}_t := \rvf(\rvx_t, t) - \tilde{\rvf}(\rvx_t, t)$ represent the difference in drifts for the given path $\rvx_t$.

\paragraph{Girsanov's Likelihood Ratio.}
Girsanov's theorem provides a fundamental likelihood ratio between these two ways of interpreting \emph{the same observed path}. It states:
\begin{mdframed}
\[
\frac{p_{\rvf}(\rvx_{0:T})}{p_{\tilde{\rvf}}(\rvx_{0:T})} = \exp\left( \int_0^T \bm{\delta}_t^\top g(t)^{-1} \diff\rvw_t - \frac{1}{2} \int_0^T \| g(t)^{-1} \bm{\delta}_t \|^2 \diff t \right).
\]
\end{mdframed}
This compact formula is an exponential of two integrals. The first is an It\^o integral, while the second is a standard Riemann integral. This ratio is crucial in diffusion models, allowing us to bridge between different data generation processes and to evaluate model likelihoods.

Girsanov's theorem is best understood as a change-of-variable formula for \emph{measures}. Just as a change of variables in calculus transforms an integral between coordinate systems via the Jacobian determinant, Girsanov's theorem provides the corresponding factor (the Radon–Nikodym derivative) to transform probabilities or expectations between two stochastic processes, when the drift changes but the diffusion remains the same.




\subsection{Girsanov’s Theorem as a Bridge Between Likelihood Training and Score Matching}\label{subsec:girsanov}

After understanding how Girsanov's theorem relates the likelihoods of a single
path under different drift assumptions, we now delve into its implications for
diffusion models.

Recall the forward SDE in diffusion models:
\[
\diff \rvx_t = \rvf(\rvx_t, t) \diff t + g(t) \diff \rvw_t,
\]
which induces a \emph{path distribution} $P$ over full trajectories
$\rvx_{0:T} := \{\rvx_t\}_{t=0}^T$ (that is, the joint law of the process over
the entire time interval).
The reverse-time SDE, parameterized by a learnable score function
$\rvs_{\bm{\phi}}(\rvx_t, t)$, is given by
\[
\diff \rvx_t = \big[\rvf(\rvx_t, t) - g^2(t) \rvs_{\bm{\phi}}(\rvx_t, t)\big] \diff t
+ g(t) \diff \bar{\rvw}_t,
\]
which in turn defines another path distribution
$P_{\bm{\phi}}$ over trajectories.


\paragraph{Two Notions in Diffusion Models.}
In diffusion models, we navigate between two core perspectives for describing the stochastic process $\rvx_{0:T}$: the forward process and its reverse-time counterpart. These perspectives give rise to two distinct but related objectives:
\begin{itemize}
    \item \textbfs{Concept 1. Marginal Distribution Matching}: This goal constructs a reverse-time process whose marginals $p_t(\rvx_t)$ match those of the forward SDE, starting from noise at time $T$ and recovering the data distribution at $t=0$. As emphasized, the Fokker–Planck equation ensures this marginal consistency for the reverse-time SDE.
    \item \textbfs{Concept 2. Joint Path Distribution Matching}: This stronger objective seeks to match the full joint distribution over the entire trajectory $P = p(\rvx_{0:T})$. Rather than just matching snapshots at individual time steps, this condition ensures that the entire sequence of states and their temporal dependencies are faithfully reproduced.
\end{itemize}

Matching the full path distribution $P$ ensures all marginals match. Formally, let $\rvx_{0:T} := \{\rvx_t|t \in [0,T]\}$ be a stochastic process with joint distribution $p(\rvx_{0:T})$. Suppose another process with joint $q(\rvx_{0:T})$ satisfies
$$ p(\rvx_{0:T}) = q(\rvx_{0:T}). $$
Then for any $t \in [0,T]$, the marginal distributions are
\begin{align*}
p_t(\rvx_t) = \int p(\rvx_{0:T})  \diff\rvx_{[0,T]\setminus \{t\}}, \quad 
q_t(\rvx_t) = \int q(\rvx_{0:T})  \diff\rvx_{[0,T]\setminus \{t\}},
\end{align*}
which implies
$$ p_t(\rvx_t) = q_t(\rvx_t), \quad \forall t \in [0,T]. $$
Thus, joint path matching implies marginal matching.

However, the reverse is not true: two processes may share identical marginals at every time step yet differ significantly in their temporal correlations. Marginal matching lacks the ability to capture these inter-time dependencies, which are encoded only in the joint distribution.


\paragraph{Girsanov Bridges the Two Goals.}
While reverse-time SDEs are primarily designed for marginal matching (Concept 1), Girsanov's theorem reveals a deeper connection: score matching across time also encourages joint path matching (Concept 2).

More precisely, Girsanov’s theorem relates the forward path distribution $P$ and the learned reverse path distribution $P_{\bm{\phi}}$. The objective function in score-based diffusion models is the KL divergence between these path measures:
\begin{align}\label{eq:girsanov-KL} \mathcal{D}_{\mathrm{KL}}(P  \|  P_{\bm{\phi}}) =  \frac{1}{2} \mathbb{E}_{P} \left[\int_{0}^{T} g^2(t) \big\|\rvs_{\bm{\phi}}(\rvx_t,t) - \nabla_{\rvx}\log p_t(\rvx_t)\big\|^2 \diff t\right] + \mathrm{Const.}, \end{align}
Here, the constant does not depend on $\bm{\phi}$, and we use the fact that the It\^o integral has zero expectation under $P$. This expression shows that minimizing KL divergence between joint paths is equivalent to learning a score function $\rvs_{\bm{\phi}}$ that approximates the true score $\nabla_{\rvx_t} \log p_t(\rvx_t)$. Thus, score matching, although framed as a marginal objective, effectively promotes alignment of the entire joint path distribution.

\paragraph{Implicit Likelihood Training.}
Beyond just matching path distributions, score matching implicitly allows diffusion models to achieve a fundamental goal of generative modeling: approximating the data likelihood~\citep{song2021maximum}. 

The connection becomes clear through a powerful concept called the change-of-measure formula. This formula, illuminated by Girsanov's theorem, allows us to express the logarithm of the marginal likelihood of the data at $t=0$ ($p_{\bm{\phi}}(\rvx_0)$) under our learned model.
\begin{align}\label{eq:girsanov-likelihood}
    \log p_{\bm{\phi}}(\rvx_0) = \log \int p_T(\rvx_T) \cdot \frac{p_{\bm{\phi}}(\rvx_{0:T})}{p(\rvx_{0:T})} p(\rvx_{0:T}) \diff \rvx_{0:T}.
\end{align}
Here, $p_T(\rvx_T)$ is the known distribution of noise at time $T$ given by the forward SDE. The term $\frac{p_{\bm{\phi}}(\rvx_{0:T})}{p(\rvx_{0:T})}$ is the density ratio between the learned reverse process and the forward process for a given path $\rvx_{0:T}$—an object precisely quantified by Girsanov's theorem. Essentially, this formula calculates the likelihood of generated data by re-weighting the known likelihood of noise based on how well our learned reverse dynamics explain the observed path's trajectory.

We further draw connection back to the KL minimization in \Cref{eq:girsanov-KL} which concerns the discrepancy between the full forward path distribution and the learned reverse path distribution. The two, \Cref{eq:girsanov-KL} and \Cref{eq:girsanov-likelihood} are deeply intertwined: optimizing this score matching objective (the training loss) directly translates to learning the Girsanov density ratio, thereby implicitly maximizing the data likelihood ($p_{\bm{\phi}}(\rvx_0)$). This elegant connection beautifully ties together Girsanov’s theorem, score-based learning, and the ultimate generative modeling goal of assigning high probability to real data.
